\section{Qualitative Use Before Calibration}

The framework can guide investigation before its coefficients are known. Classify constraints using public criteria and record the evidence supporting each classification.

\subsection{Eight Regions}

For a coarse view, mark each coordinate low or high under a declared scale. The three binary classifications yield eight regions.

\begin{table}[ht]
\centering\small
\begin{tabularx}{\textwidth}{@{}cccYY@{}}
\toprule
$R$ & $T$ & $G$ & Initial focus & Questions to investigate\\
\midrule
Low & Low & Low & The remaining service boundary & Is demand, quality, or an omitted constraint limiting use?\\
High & Low & Low & Resource barrier & Can substitution, efficiency, or scale improve access?\\
Low & High & Low & Time & Which waiting periods supply information or complete a process?\\
Low & Low & High & Coupling & Can interfaces or dependencies be isolated?\\
High & High & Low & Resources and time & Which stages can be redesigned or parallelized?\\
High & Low & High & Resources and coupling & Can a different architecture reduce both?\\
Low & High & High & Time and coupling & Which interacting responses determine the schedule?\\
High & High & High & The full conversion chain & What phased or alternative route is feasible?\\
\bottomrule
\end{tabularx}
\caption{Regions identify an investigation priority. They do not rank total cost or predict a launch date.}
\end{table}

Near a classification boundary, report the neighboring categories and the evidence that could distinguish them. The cost of coarse classification is loss of within-category detail.

\subsection{Comparison Without a Full Weight Vector}

For comparable states under the exponential model,
\begin{equation}
 \ln\frac{\eta_a}{\eta_b}
 =\alpha\Delta R+\beta\Delta T+\gamma\Delta G.
\end{equation}
With positive coefficients and matching scales, a state no higher on any constraint and lower on at least one has a lower modeled response. This dominance relation does not require exact weights.

Mixed changes need more information. A reduction in $R$ and an increase in $G$ cannot be ranked from their signs alone. Coefficient ratios, magnitudes, and uncertainty bounds determine the result.

Ordinal grades support within-coordinate order. They do not establish equal intervals or a common scale across coordinates. Grade “four” in resources cannot be compared directly with grade “three” in time.

\subsection{Direction Over Time}

With fixed baseline and coefficients,
\begin{equation}
 \frac{\mathrm d\ln\eta}{\mathrm dt}
 =\alpha\dot R+\beta\dot T+\gamma\dot G.
\end{equation}
Common-direction changes permit a sign judgment under positive weights. Mixed directions require bounds or estimates. If coefficients and baseline change, their derivative terms must be included.

This is enough to organize some useful questions: is the physical constraint falling while review remains unchanged? Does a new interface reduce resource use but expand dependencies? Is a reported gain a change in the task boundary?

The qualitative account should end with evidence needs. It can identify an unresolved tradeoff without inventing a numerical answer.
